\subsubsection{Baseline Reproduction and Implementation Scope}
\label{subsubsec:baseline_reproduction}

The reference study models traffic as an average per-node load, whereas
our simulator generates one network-wide request stream. Its Table~IV
was therefore used as the low-load reproduction check when fixing the
topology-specific load mappings, which were then held constant across
all schemes and loads. Supplementary Table~S1 compares the published
values with our reproduced ten-seed means.

As a secondary multi-load check, Supplementary Table~S2 compares the
published values in Table~V of the reference study with our reproduced
values. Because the reference study does not state its aggregation formula,
this source-style check transparently averages the five load-specific signed
percentage changes. This calculation is distinct from
Eqs.~\eqref{eq:load_averaged_metric}--%
\eqref{eq:load_averaged_fsu_change}, which first average the metric
values over \(\rho\) and are used for the principal \acs{aca}
comparisons. For K1-RSA, all twelve topology--metric differences are
within \(1.33\) percentage points of the published values.

The reference study does not specify how its three fixed K3-RSA routes
are generated or ordered. Our implementation precomputes the three
shortest physical-distance paths and attempts them sequentially using
First-Fit. Because this choice cannot be matched unambiguously, K3-RSA
is treated as a documented common-simulator implementation rather than
an exact reproduction. Its larger TI differences are disclosed in the
supplement and it is excluded from mechanism-attribution claims.

